\title{Speculative Decoding: Exploiting Speculative Execution for Accelerating Seq2seq Generation}

\begin{abstract}
We propose Speculative Decoding (SpecDec), for the first time ever\footnote{This work was initially announced in \textbf{March 2022} (\url{https://arxiv.org/abs/2203.16487}) under the name \textit{Generalized Aggressive Decoding}. It has been formally renamed \textit{Speculative Decoding} in our submission to ICLR'23, which has been publicly available since \textbf{September 2022} at \url{https://openreview.net/pdf?id=H-VlwsYvVi}. This marks \textbf{the first time} "Speculative Decoding", which explicitly studies the idea of speculative execution to accelerate Transformer inference, has been proposed.}, to formally study exploiting the idea of speculative execution to accelerate autoregressive (AR) decoding. Speculative Decoding has two innovations: Spec-Drafter -- an independent model specially optimized for efficient and accurate drafting -- and Spec-Verification -- a reliable method for verifying the drafted tokens efficiently in the decoding paradigm. Experimental results on various seq2seq tasks including machine translation and abstractive summarization show our approach can achieve around $5\times$ speedup for the popular Transformer architectures with comparable generation quality to beam search decoding, refreshing the impression that the \textit{\textit{draft-then-verify}} paradigm introduces only $1.4\times$$\sim$$2\times$ speedup. In addition to the remarkable speedup, we also demonstrate 3 additional advantages of SpecDec, revealing its practical value for accelerating generative models in real-world applications. Our models and codes are available at \url{https://github.com/hemingkx/SpecDec}.
\end{abstract}

\section{Introduction}\label{sec:intro}
As the \textit{de facto} method for text generation, AutoRegressive (AR) decoding is widely blamed for its poor inference efficiency due to its low level of parallelism, which fails to utilize the full potential of modern parallel computing devices like GPUs. This inefficiency not only leads to high deployment costs but also limits the application of advanced AR models in real-time scenarios.

In this work, we study the \textit{draft-then-verify} paradigm for accelerating seq2seq generation of an existing AR model\footnote{The existing AR model in this paper refers to the targeted Transformer using AR decoding that we want to accelerate.}. As shown in Figure~\ref{fig:draft-then-verify}, the \textit{draft-then-verify} paradigm first generates a number of drafted tokens efficiently and then verifies these tokens using the existing AR model in parallel to ensure the decoding result matches AR decoding. However, previous attempts in the ``\textit{draft-then-verify}'' paradigm such as Blockwise Decoding~\citep{Stern:2018blockwise} and Aggressive Decoding~\citep{Sun:2020SAD} tend to lack in-depth investigation of this paradigm. Their modest speedup (i.e., $1.4\times$$\sim$$2.0\times$) or limitation to certain seq2seq tasks like Grammatical Error Correction (GEC) has caused this paradigm to be underestimated, resulting in it not receiving much attention and remaining dormant for years.

To fully exploit the \textit{draft-then-verify} paradigm, we propose Speculative Decoding (SpecDec), drawing inspiration from speculative execution\footnote{Speculative execution is an optimization technique used in computer architecture where a system performs some task in advance to avoid delays that would have to be incurred by doing the task after it is known that it is required (\url{https://wikipedia.org/wiki/Speculative_execution}).} in computer architecture, with two key innovations that improve drafting and verification processes respectively. For drafting, we derive two principles for designing the drafting model\footnote{The drafting model is also called the drafter in this paper.}: \textit{the Capability Principle} and \textit{the Latency Principle}. Following these two principles, we propose Spec-Drafter -- a specialized independent model optimized in the \textit{draft-then-verify} paradigm, which can accurately and efficiently fulfill the drafting task.

For verification, we propose an advanced method -- Spec-Verification that relaxes the vanilla verification strategy. Spec-Verification allows the decoding results of SpecDec to be slightly different from AR greedy decoding, offering an opportunity to accept more drafted tokens without sacrificing generation quality and leading to higher decoding efficiency. 

We conduct extensive experiments on various seq2seq generation tasks like machine translation and abstractive summarization. Results show our approach can achieve around $5\times$ speedup for the popular Transformer architectures with comparable generation quality to beam search decoding, largely outperforming previous \textit{draft-then-verify} work ($1.4\times$$\sim$$2.0\times$ speedup). Moreover, we demonstrate that SpecDec has several additional advantages that enhance its practicality for accelerating generative models in real-world applications.

Our contributions can be summarized as follows:

\begin{itemize}
    \item We are the first work that explicitly exploits the idea of speculative execution to accelerate Transformer inference. Our proposed two key innovations -- the independent Spec-Drafter and Spec-Verification strategy allow SpecDec to achieve over $5\times$ lossless speedup over autoregressive decoding in seq2seq tasks, refreshing the impression that the ``\textit{draft-then-verify}'' paradigm only has a limited $1.5\times \sim 2\times$ acceleration potential.
    \item We demonstrate 3 advantages of SpecDec with extensive empirical results in addition to its remarkable acceleration performance: better latency-throughput trade-off, easy adaptability for existing models and retaining the behavior of the original model, revealing its huge practical value and bringing the long-dormant \textit{draft-then-verify} paradigm back into the spotlight.
\end{itemize}

\section{Background: \textit{draft-then-verify} decoding}
\label{sec:Background}

The ``\textit{draft-then-verify}'' paradigm first \textit{drafts} multiple tokens efficiently, as a \textit{speculation} of AR decoding results; then, it \textit{verifies} these tokens in parallel to ensure they match the AR decoding result, as illustrated in Figure \ref{fig:draft-then-verify}. It is an implicit implementation of speculative execution in Transformer inference.

\paragraph{Draft} There are different approaches to drafting tokens, including model-based \cite{Stern:2018blockwise} and input-(context-) based\footnote{This kind of method is usually limited to special tasks like GEC and retrieval-augmented generation. In this paper, we mainly focus on the model-based methods.} methods~\cite{Sun:2020SAD,yang2023inference}. Take Blockwise Decoding -- the most representative work attempting the \textit{draft-then-verify} paradigm -- as an example (illustrated in Figure \ref{fig:NA-Block}(a)): it introduces additional $k-1$ feedforward network (FFN) heads on top of an \textit{existing} AR model, enabling the model to predict the next $k$ drafted tokens in parallel during inference.

\paragraph{Verify} The generated drafted tokens are fed into the original AR model and verified in parallel. Specifically, it finds the bifurcation position $c$, the largest index that ensures all previous $c-1$ drafted tokens and the corresponding AR decoded tokens are identical:

\begin{equation}\small
c = \arg \max_i \frac{\mathbb{I}(\widetilde{y}_{j+i} \neq \hat{y}_{j+i})}{i}, 1 \le i \le k 
\end{equation}

\begin{equation}\label{eq:strict-verify}\small
\hat{y}_{j+i} = \arg \max_{y} \log P(y\mid \hat{\bm{y}}_{\le j},\widetilde{\bm{y}}_{j+1 \cdots j+i-1}, \bm{x}; \bm{\theta}_{\textrm{AR}})
\end{equation}

where $\mathbb{I}(\cdot)$ is the indicator function, $\bm{x}$ is the source sentence, $\hat{\bm{y}}_{\le j}$ is the previously generated tokens\footnote{We use $\widetilde{\bm{y}}$ to denote drafted tokens, while we use $\hat{\bm{y}}$ to denote AR decoded/verified generation results.} and $\widetilde{y}_{j+i}$ is the $i$-th drafted token. Drafted tokens after the position $c$ are all discarded. The final decoded tokens in the current iteration are:

\begin{equation}
\small
    \hat{\bm{y}}_{j+1 \cdots j+c} = (\widetilde{\bm{y}}_{j+1  \cdots  j + c - 1}, \hat{y}_{j+c})
\end{equation}

The above draft and verification steps are iterated until the termination condition is met. 

\section{Speculative Decoding}
\label{sec:SpecDec}
To fully exploit speculative execution for Transformer inference, we propose Speculative Decoding (SpecDec) with two innovations -- Spec-Drafter and Spec-Verification that substantially improve drafting (Section 3.1) and verification (Section 3.2) respectively.

\subsection{Spec-Drafter}
\subsubsection{Design Principles}

As a crucial ingredient in the \textit{draft-then-verify} paradigm, the drafting process has a drastic impact on end-to-end acceleration performance. However, there are very limited explorations of the designing principles for the drafter by previous studies -- most of them arbitrarily implement a drafter, which accounts for their undesirable acceleration results.

To understand the effect of drafting, we look into the overall latency in the \textit{draft-then-verify} paradigm for one sample of the length $L$ as follows:

\begin{equation}\label{eq:latency}\small
T = \underbrace{\frac{L}{\textrm{\bf Tok.}}\times t_d}_{{\footnotesize\raggedright \textrm{total drafting latency}}} + \underbrace{\frac{L}{\textrm{\bf Tok.}}\times t_v}_{{\footnotesize\raggedright \textrm{total verification latency}}}
\end{equation}

where \textbf{Tok.} denotes the average number of drafted tokens accepted per iteration, $t_d$ and $t_v$ are the time costs of drafting and verification\footnote{We don't discuss $t_v$ in this paper because it is determined by the existing AR model and thus regarded constant.} each iteration respectively.

According to Eq (\ref{eq:latency}), \textbf{Tok.} is inversely proportional to the number of iterations, which is primarily influenced by drafting accuracy: A drafter that is more capable of drafting can attain greater \textbf{Tok.} values, consequently completing the decoding process in fewer iterations. This observation leads us to derive the first principle for designing the drafter:

\begin{quote}
    \textbf{Principle I (Capability Principle)}: The drafter model should be seriously invested to guarantee its capability of accurate drafting.
\end{quote}

Principle I is the most crucial principle in determining the end-to-end speedup, as it directly influences the value of \textbf{Tok.} which affects both total drafting and verification latency. Surprisingly, little previous work adheres to this seemingly simple and straightforward principle maybe due to the concern of increasing the drafting latency. For instance, the drafter in Blockwise Decoding is not properly invested: Its drafter not only has limited parameters (FFN prediction heads), making it difficult to fit the challenging drafting task, but more importantly, it employs a shared attention mechanism that forces all drafted tokens to share a single set of attentions (only differentiating at the final prediction head), as shown in Figure~\ref{fig:NA-Block}(a). However, different target positions should attend different context tokens, as illustrated in Figure~\ref{fig:heatmap}. Despite its computational efficiency, the shared attention mechanism in Blockwise decoding severely constrains the drafter's capability, resulting in low drafting accuracy and consequently leading to most drafted tokens being discarded.  

In addition to the drafter's accuracy, its latency also impacts the end-to-end speedup result but from another perspective -- by affecting the latency of each iteration (i.e., $t_d$ in Eq~(\ref{eq:latency})) -- from which we derive Principle II for designing the drafter:

\begin{quote}
    \textbf{Principle II (Latency Principle)}: The drafter should be fast at generating drafted tokens to minimize the latency overhead of each iteration.
\end{quote}

Designing a fast drafter solely based on Principle II is not difficult, as done in most previous work. The real challenge lies in designing a low-latency drafter without compromising its capability (Principle I), since it is difficult to achieve both low latency and high capability simultaneously.

\subsubsection{Model Architecture} We propose Spec-Drafter, which adheres to both principles for accurate and fast drafting. To ensure the drafter is sufficiently capable of accurate drafting (Principle I), Spec-Drafter employs an independent encoder-decoder model architecture, which generates drafted tokens conditioned on the leftward context and source tokens in a mask-predict manner~\cite{Ghazvininejad:2019}, as illustrated in Figure~\ref{fig:NA-Block}(b). This independent model design facilitates Spec-Drafter to predict each drafted token using distinct attention queries, in contrast to Blockwise Decoding employing a shared attention query for predicting all drafted tokens (as illustrated in Figure~\ref{fig:NA-Block}). In this way, Spec-Drafter could better align with the AR model's behavior, thereby increasing the chances of its drafted tokens being accepted during verification, as shown in Figure~\ref{fig:heatmap}. 

To make Spec-Drafter fast (Principle II) without compromising its capability, we design its decoder to be lightweight by reducing the number of decoder layers and reallocating the freed-up budget to its encoder (by increasing its depth), which is motivated by the fact that the encoder is forwarded only once, while the decoder is frequently forwarded for iterative decoding. This encoder-favored modeling has been demonstrated by previous work to improve latency with little generation quality degradation~\cite{Kasai:2021desd, Sun:2020SAD, Ge:2022}. We find it also highly effective for the drafter in the \textit{draft-then-verify} decoding paradigm.

\subsubsection{Training and Inference} Formally, given the source sentence $\bm{x}$ and the randomly sampled prefix $\bm{y}_{\le p}$ ($0 \le p < m$)  of the target sentence, Spec-Drafter appends $k$ special ``\texttt{[MASK]}'' tokens to $\bm{y}_{\le p}$, and is trained to predict these masked tokens in parallel:

\begin{small}
\begin{equation}\label{eq:nat-drafter}
\mathcal{L}_{\textrm{Spec-Drafter}}=\sum_{i=p+1}^{p+k} \log P\left(y_{i} \mid \bm{y}_{\le p}^k, \bm{x}; \bm{\theta}_{\textrm{Spec-Drafter}}\right)
\end{equation}
\end{small}

\begin{small}
\begin{equation}
\bm{y}_{\le p}^k=(y_1,\cdots,y_p, \underbrace{[\rm{MASK}], \cdots, [\rm{MASK}]}_{\times k})
\end{equation}
\end{small}

In addition, we leverage the glancing strategy following \citet{Qian:2020}, which exploits curriculum learning during training to get better generation performance. 

During inference, Spec-Drafter appends $k$ ``\texttt{[MASK]}'' tokens to the previously decoded tokens $\hat{\bm{y}}_{\le j}$ and simultaneously predict these masked tokens as a drafted block:

\begin{equation}\label{eq:na-draft-infer}\small
\widetilde{y}_{j+i} = \arg \max _{y} \log P\left(y\mid \hat{\bm{y}}_{\le j}^k, \bm{x};\bm{\theta}_{\textrm{Spec-Drafter}}\right)
\end{equation}

where $i=1,\ldots,k$.

\subsection{Spec-Verification}
\label{sec:SpecDec++}
As introduced in Section~\ref{sec:Background}, the vanilla verification strategy of preliminary studies only accepts the drafted tokens that match the top-1 result of the AR model, which guarantees that the decoding results are identical to AR greedy decoding. However, the top-1 results are not necessarily better than the drafted tokens, especially when the paradigm is equipped with a high-quality drafter. Therefore, the strict verification criterion (i.e., top-1 matching) will result in many good drafted tokens being discarded just because they are different from the top-1 result of the AR model, which limits the speedup of the paradigm.

To make better use of the drafting results, we propose an advanced verification strategy named Spec-Verification, which is illustrated in Figure~\ref{fig:SpecDec++}. Instead of the rigid matching requirement shown in Eq~(\ref{eq:strict-verify}), Spec-Verification relaxes the criterion to trust the drafting results more, by only requiring the drafted tokens to fall in \textit{top-$\beta$} candidates with a tolerable (log-likelihood) score gap $\tau$ (away from the top-1 result). Formally, it will accept the $i$-th drafted token $\widetilde{y}_{j+i}$ if all previous $i-1$ tokens are accepted, and Eq~(\ref{eq:SpecDec++1}) and~(\ref{eq:SpecDec++2}) are both true:

\begin{small}
\begin{equation}\label{eq:SpecDec++1}
   \log P(\widetilde{y}_{j+i}|\triangle;\bm{\theta}_{\textrm{AR}}) \ge \log P(\hat{y}_{j+i}^{(\beta)}|\triangle;\bm{\theta}_{\textrm{AR}}),
\end{equation}
\end{small}

\begin{small}
\begin{equation}\label{eq:SpecDec++2}
    \log P(\hat{y}_{j+i}^{(1)}|\triangle;\bm{\theta}_{\textrm{AR}}) - \log P(\widetilde{y}_{j+i}|\triangle;\bm{\theta}_{\textrm{AR}}) \le \tau,
\end{equation}
\end{small}

\begin{small}
\begin{equation}
    \triangle = \hat{\bm{y}}_{\le j}, \widetilde{\bm{y}}_{j+1 \cdots j+i-1}, \bm{x},
\end{equation}
\end{small}

where $\log P(\hat{y}_{j+i}^{(\beta)}|\triangle;\bm{\theta}_{\textrm{AR}})$ is the top-$\beta$ ranked result's log-likelihood score by the AR model.

\section{Experiments}
\subsection{Experimental Settings}
\label{sec:settings}
\paragraph{Datasets and Evaluation} We mainly evaluate our approach on two standard machine translation benchmarks: WMT14 EN$\leftrightarrow$DE (4.5M pairs) and WMT16 EN$\leftrightarrow$RO (610K pairs). Following prior work~\citep{ott-etal-2018-scaling}, for WMT14 EN$\leftrightarrow$DE translation, we adopt \textit{newstest-13} as our validation set for finding the best hyperparameters, and test on \textit{newstest-14}. For WMT16 EN$\leftrightarrow$RO translation, we use the dataset released by \citet{Lee:2018}, where \textit{newsdev2016} and \textit{newstest2016} are taken as validation and test sets. We use 32K Byte Pair Encoding (BPE)~\citep{Sennrich:2016} subwords\footnote{We use the same BPE tokenization and vocabulary as ~\citet{Ghazvininejad:2019}.} as the joint source-target dictionary. We evaluate performance with BLEU~\citep{Papineni:2002} for both language pairs\footnote{We also report sacreBLEU \cite{Post:2018} and COMET \cite{Rei:2020} scores in Appendix \ref{sec:other-evaluation}.}. 

For inference efficiency, we report decoding speedup over beam search. Specifically, we test the inference speed by running the model with one sentence at a time (batch=1). We perform model inference with fairseq implementation\footnote{\url{https://github.com/pytorch/fairseq}} using Pytorch 1.10.1 with 1 Nvidia Tesla P100-PCIe of 16GB GPU memory under CUDA 11.1.

\paragraph{Model Configuration} The primary target model we accelerate in our experiments is the Transformer-base model with a 6-layer encoder and a 6-layer decoder of 512/2048 embedding/FFN dimension, which can achieve state-of-the-art results on the benchmarks under comparable model size conditions. For the Spec-Drafter, we adopt a similar architecture to the AR model except with 12 encoder layers and 2 decoder layers to make sure it adheres to both the Capability and Latency principles. We apply sequence-level knowledge distillation~\cite{kim2016sequence} by the AR teacher to the Spec-Drafter to align its behavior with the AR model as much as possible. We include model training details in Appendix \ref{sec:hyperparameters}. For the Spec-Verification, we find the hyperparameters $\beta$ and $\tau$ leading to the best generation quality on the validation set. Besides, we re-implement Blockwise Decoding\footnote{In the original paper of Blockwise Decoding, there is also a variation that allows the AR model to be fine-tuned for better drafting tokens. We don't discuss this variation because it severely affects the generation quality.} using the same device and environment as ours to facilitate fair comparison.

\subsection{Results}\label{subsec: result} We present the performance and the acceleration effect of SpecDec to Transformer in Table~\ref{tab:exp}. As reported in the previous work~\cite{Stern:2018blockwise}, Blockwise Decoding ($k=10$) can only achieve $1.4\times$$\sim$$2\times$ speedup without affecting the generation results over the Transformer-base model. Further increasing the parallel capability of Blockwise Decoding (e.g., $k=25$) will not introduce more speedup as its limited drafting accuracy prevents more drafted tokens from being accepted. In contrast, our SpecDec shows consistent performance improvement with increased parallel capabilities ($k=10$ $\rightarrow$ $k=25$), resulting in around $4.6\times$$\sim$$5.5\times$ speedup across the translation benchmarks; moreover, it even achieves an improvement in generation quality (by the BLEU metric) compared with AR greedy decoding. 

Similar results are also observed when accelerating the Transformer with a 12-layer encoder and 2-layer decoder -- SpecDec can still achieve around $2.5\times$$\sim$$3.3\times$ speedup while Blockwise Decoding's acceleration effect becomes nearly negligible over the fast AR baseline.

\subsection{Analysis}
In this section, we conduct a comprehensive and thorough analysis, to demonstrate that the significant improvement of the SpecDec arises from both the Spec-Drafter (Section~\ref{subsubsec:draft_analysis}) and Spec-Verification (Section~\ref{subsubsec:verify_analysis}).

\subsubsection{Drafting}\label{subsubsec:draft_analysis}
According to Table~\ref{tab:ablation-draft}, the Spec-Drafter significantly outperforms the head-based drafter (as used in Blockwise Decoding) in terms of both end-to-end generation quality and efficiency. To further investigate the Spec-Drafter, we conduct an ablation study on the principles it follows. Ablating the Capability Principle by reducing its size results in a drastic drop in end-to-end acceleration performance, as more iterations are needed to complete the decoding process, indicated by a lower \textbf{Tok.}. When we ablate the Latency Principle by using a balanced (6+6) encoder-decoder architecture for the drafter, it also experiences a substantial decline in end-to-end acceleration performance due to increased latency in each iteration (reflected by a higher $t_d$).    

Moreover, we analyze the block size $k$'s effect on the end-to-end acceleration performance of the paradigm in Table~\ref{tab:block-size}. In contrast to the blockwise decoding achieving its best acceleration performance at $k=10$ as shown in Table~\ref{tab:exp}, SpecDec achieves its best performance at $k=25$ with 7.89 mean accepted tokens each iteration.  Further increasing $k$ has an adverse effect, because it will become very hard for the model to learn to draft too many tokens simultaneously given the model capacity, resulting in a drop of \textbf{Tok.}. 

\subsubsection{Verification}\label{subsubsec:verify_analysis}
We study the effect of Spec-Verification on the development set (i.e., \textit{newstest-2013}) of WMT14 EN$\rightarrow$DE in Table~\ref{tab:verification}. Moderately increasing $\tau$ and $\beta$ in the Spec-Verification not only leads to an increase of mean accepted tokens (\textbf{Tok.}) and speed since AR verification becomes less strict but also improves the generation quality over greedy decoding. However, the generation quality may decrease if over-relaxed: the BLEU score will degrade from the peak of 26.97 to 26.58 when decoding with \textit{top-5} selection (i.e., $\beta=5$) and $\tau=5.0$. Based on the results in the development set, we select $\beta=3, \tau=1.0$ as our Spec-Verification hyperparameters.

\subsection{Practical Value}
\label{sec:practical}
In addition to the remarkable speedup results, we demonstrate SpecDec's additional advantages that enhance its practical value in the following three aspects:

\paragraph{Better latency-throughput trade-off} SpecDec achieves inference acceleration by increasing the GPU computing parallelism. Although increasing the batch size can also increase the computing parallelism to improve throughput, it results in increased latency, which is not desirable in real-world application scenarios. Therefore, a smaller batch size is often employed during inference, but this in turn results in the underutilization of GPU computing resources, leading to the dilemma of low throughput for small batches and high latency for large batches, as illustrated by Figure~\ref{fig:latency-throughput}. SpecDec effectively addresses this dilemma. Even by maintaining a small batch size, SpecDec can fully utilize the computing performance of the GPU, significantly improving both efficiency and throughput.

\paragraph{Easily adaptable for existing models} In many practical applications, generative models are often pretrained with massive data, which exhibits very high performance. Developing a faster model from scratch to replace the pretrained model is highly challenging and typically requires substantial computational costs to reiterate the pretraining process. Otherwise, the quality of the new models is very likely to be compromised despite the increased speed, as Table~\ref{tab:sum} shows. However, SpecDec can be easily adapted to accelerate existing pretrained models. Taking the BART-base~\cite{Lewis:2020} model for the abstractive summarization task as an example, we can easily achieve $5.1\times$ speedup without compromising the generation quality only by initializing the Spec-Drafter with the BART-base encoder and training it with the BART-base distilled summarization training set. 

\paragraph{Retaining the behavior of the original model} As introduced in Section~\ref{sec:intro}, one significant advantage of SpecDec is that it does not develop a new faster model to replace the existing model. Instead, it accelerates the existing model with minimal changes to its behavior. As shown in Table~\ref{tab:relative-bleu}, the consistency (BLEU) of SpecDec's generated results with the original model exceeds 85\%, while that of a newly built fast NAR model is only around 55\%. The characteristic of maintaining the behavior of the original model makes SpecDec even more valuable in practical applications because transitioning from a well-tested and mature model to one with substantially different behavior is risky, requiring extensive recalibration and various offline evaluations and online feedback in practice. 

\section{Related Work}\label{sec:related}

\paragraph{Speculative Decoding} We have demonstrated since early 2022 (see our arXiv preprints in 2022) that our proposed methodology, which formally introduces an independent model as a drafter combined with an advanced verification strategy to fully exploit speculative execution, is promising and has potential to evolve into a \textit{de facto} standard in the future for efficient and lossless decoding. Since this work was proposed, we are pleased to see an increasing number of following studies~\cite{Leviathan:2023specdec, Chen:2023specsampling, Kim:2023bild, Spector:2023stagedspec, Zhang:2023draftverify} acknowledge, explore and adopt this methodology to accelerate Transformer inference. Among them, \citet{Leviathan:2023specdec} use the same name as ours (i.e., Speculative Decoding), employing a small AR model as a drafter\footnote{We provide a detailed comparison between \citet{Leviathan:2023specdec} and our work in Appendix~\ref{sec:comparison-with-other-work}.} as well as advanced sampling algorithm. \citet{Chen:2023specsampling} is similar to \citet{Leviathan:2023specdec} but it was the first to validate this methodology to accelerate a large language model (i.e., 70B Chinchilla) with a 4B drafter model, thus receiving the most attention. SpecInfer~\cite{Miao:2023specinfer} proposed to utilize various boost-tuned small language models for joint drafting, to improve the speculation accuracy of the LLM’s outputs. Besides, it introduces an advanced token tree verification strategy to verify all candidate token sequences in parallel. DistillSpec~\cite{Zhou:2023distillspec} further investigated the efficacy of knowledge distillation in enhancing the alignment between the target model and the drafter in speculative decoding. In addition to employing additional models as drafters, there has also been some research that proposes various strategies to efficiently generate drafts from the LLM itself~\cite{Santilli:2023paralleldecoding, Zhang:2023draftverify}. All the following research strongly backs up the value of this original work.

\paragraph{Early \textit{Draft-then-verify} attempts} This work is a generalized version of our previously proposed (Input-guided) Aggressive Decoding\footnote{As our technical report \citep{ge2022lossless} in May 2022 discusses, the Input-guided Aggressive Decoding is indeed a special case of Speculative Decoding.} \citep{Sun:2020SAD} in Grammatical Error Correction (GEC), which assumes that the input is exactly the sentence to be generated in the future and then verifies the whole sentence in parallel. Blockwise Decoding~\citep{Stern:2018blockwise} inserted $k-1$ feedforward heads on top of the Transformer decoder to generate $k$ positions in parallel and used the original head to verify these outputs. However, both the above studies did not fully investigate the potential of this paradigm and thus failed to uncover its great value for efficient seq2seq generation: \citet{Sun:2020SAD} only works for tasks whose inputs and outputs are highly similar (e.g., GEC). \citet{Stern:2018blockwise} overlooked the importance of drafting accuracy; as a result, their underinvested prediction heads severely limit the acceleration results. In contrast, we conduct thorough investigations and fully exploit speculative execution, refreshing the impression of its limited acceleration potential and revealing its real value in practice.

\paragraph{Non-autoregressive Decoding} There is also another line of work named Non-Autoregressive Decoding (NAR)~\cite{gu:2018}, which decodes multiple tokens in parallel compared with conventional AR, thus showing remarkable superiority in inference efficiency. Recently, various attempts have been made to improve the performance of NAR models, including training with alignment-based objectives~\cite{Libovicky:2018, Ghazvininejad:2020a, Saharia:2020, Gu:2021, Shao:2022ctc}, modeling dependencies between target tokens~\cite{Ghazvininejad:2019, Shu:2020, Qian:2020, Bao:2021} and designing various model architectures~\cite{Zheng:2021, Huang:2022dat}. As discussed in Section~\ref{sec:practical}, replacing a powerful pretrained model with NAR models in practice is challenging due to the substantial computational costs required to reiterate the pretraining process. Additionally, transitioning from a well-tested and mature model to a new NAR model with significantly different behavior poses risks in practical applications. In contrast, our proposed SpecDec can be conveniently adapted to speed up \textit{existing} AR models including high-performance pretrained models like BART with little effort. Moreover, SpecDec minimally alters the behavior of \textit{existing} models, showcasing its ability to preserve reliable generation performance in real-world practical applications.

\section{Conclusion}
We present Speculative Decoding (SpecDec), the first work to explicitly embrace the idea of speculative execution for seq2seq generation acceleration with a formal study and extensive discussion of both drafting and verification phases. Contrary to the common belief that an increase in model complexity tends to hamper inference speed, SpecDec’s introduction of an appropriately invested auxiliary drafter model substantially speeds up Transformer inference, owing to higher computational parallelism introduced by speculative execution to better utilize computing resources. 

The remarkable acceleration performance, combined with the advantages demonstrated in our experiments, clearly illustrates that SpecDec is a practical acceleration method for model deployment in real-world applications. We hope that our preliminary study could draw more attention to this promising decoding paradigm that may potentially evolve into a \textit{de facto} standard for efficient Transformer decoding in the near future.

\section*{Limitations}\label{subsec:limitation} Compared with conventional autoregressive decoding, SpecDec introduces an extra Spec-Drafter module for ensuring its drafting accuracy, which brings additional memory cost at test time. Therefore, SpecDec is particularly suitable for inference scenarios where GPU memory is abundant but there is an urgent need to improve latency -- it provides a solution to trading the surplus GPU memory for speed improvements. As thoroughly discussed in Appendix~\ref{sec:gpu_memory}, such scenarios are very common in practice. Most importantly, memory is no longer the bottleneck for practical model deployment. With the emergence and maturity of various data/tensor/pipeline parallelism techniques, the addition of more GPUs can easily address memory issues, which is also why models continue to grow larger. In contrast, latency remains an inescapable bottleneck in model deployment that cannot be resolved merely by increasing the number of machines. Therefore, we believe the increased memory consumption may not severely affect its practical value.

\end{document}
